\section{Понятие меры и интеграла Лебега}
\label{sec:measure-lebesgue-integral}

\subsection{$\sigma$-алгебра и мера}

Пусть $X$ --- некоторое множество. Семейство $\mathcal F$ его подмножеств
называется \emph{$\sigma$-алгеброй}, если
\begin{enumerate}
    \item $X\in\mathcal F$;
    \item $A\in\mathcal F\Rightarrow X\setminus A\in\mathcal F$;
    \item $A_k\in\mathcal F$ для всех $k\Rightarrow
          \displaystyle\bigcup_{k=1}^{\infty}A_k\in\mathcal F$.
\end{enumerate}
Из этих свойств следует замкнутость относительно счётных пересечений и
разностей множеств. Пара $(X,\mathcal F)$ называется измеримым пространством.

\emph{Мера} --- это функция множеств
\[
    \mu:\mathcal F\to[0,+\infty],
\]
для которой $\mu(\varnothing)=0$ и выполняется счётная аддитивность: если
$A_i\cap A_j=\varnothing$ при $i\neq j$, то
\[
    \mu\!\left(\bigcup_{k=1}^{\infty}A_k\right)
    =\sum_{k=1}^{\infty}\mu(A_k).
\]
Отсюда следуют монотонность и счётная полуаддитивность меры.

Для $X=\mathbb R^n$ важнейшим примером является \emph{мера Лебега}. На
прямоугольниках она совпадает с обычным объёмом, а затем продолжается на
$\sigma$-алгебру измеримых по Лебегу множеств. Множество $A$ называется
множеством меры нуль, если $\mu(A)=0$. Свойство, нарушающееся только на таком
множестве, говорят, выполняется \emph{почти всюду}.

\subsection{Измеримые функции}

Функция $f:X\to\overline{\mathbb R}$ называется измеримой, если прообраз
борелевских множеств измерим. Для вещественной функции эквивалентно требовать,
например, чтобы для любого $a\in\mathbb R$ множество
\[
    \{x\in X:\ f(x)<a\}
\]
принадлежало $\mathcal F$.

Сумма, разность и произведение измеримых почти всюду конечных функций также
измеримы; при ненулевом знаменателе измеримо и отношение. Поточечные пределы,
верхние и нижние пределы последовательностей измеримых функций также остаются
измеримыми.

\subsection{Построение интеграла Лебега}

Идея интеграла Лебега состоит в том, что измеряется не разбиение области по
координате, как в интеграле Римана, а множества, на которых функция принимает
определённые значения.

Для простой измеримой функции
$s=\sum_{k=1}^{m}c_k\chi_{E_k}$, $c_k\geq0$, полагают
$\int s\,d\mu=\sum_{k=1}^{m}c_k\mu(E_k)$. Для ограниченной измеримой
функции на множестве конечной меры интеграл можно ввести через верхние и нижние
суммы. Для измеримого разбиения
$T=\{E_k\}_{k=1}^m$ полагают
\[
    S(T)=\sum_{k=1}^{m}M_k\mu(E_k),\qquad
    s(T)=\sum_{k=1}^{m}m_k\mu(E_k),
\]
где $M_k=\sup_{E_k}f$, $m_k=\inf_{E_k}f$. Верхний и нижний интегралы
определяются как
\[
    \overline I=\inf_T S(T),\qquad
    \underline I=\sup_T s(T).
\]
Если $\overline I=\underline I$, их общее значение называют интегралом Лебега.

Для неотрицательной измеримой функции, возможно неограниченной, интеграл
определяется предельным переходом по возрастающим ограниченным приближениям.
Для произвольной измеримой функции вводятся положительная и отрицательная части
\[
    f^+=\max(f,0),\qquad f^-=\max(-f,0),\qquad f=f^+-f^-.
\]
Если хотя бы один из интегралов $\int f^+d\mu$, $\int f^-d\mu$ конечен, можно
определить
\[
    \int f\,d\mu=\int f^+d\mu-\int f^-d\mu.
\]
Функция называется суммируемой, если
\[
    \int_X|f|\,d\mu<\infty.
\]

\subsection{Основные свойства интеграла}

Для суммируемых функций выполняются линейность и монотонность:
\[
    \int(\alpha f+\beta g)d\mu
      =\alpha\int f\,d\mu+\beta\int g\,d\mu,
\]
\[
    f\leq g\ \text{п.в.}\quad\Rightarrow\quad
    \int f\,d\mu\leq\int g\,d\mu.
\]
Кроме того,
\[
    \left|\int f\,d\mu\right|
    \leq\int|f|\,d\mu.
\]
Изменение функции на множестве меры нуль не изменяет значения интеграла.
Именно поэтому в пространствах Лебега естественно отождествлять функции,
равные почти всюду.

\subsection{Предельный переход под знаком интеграла}

\emph{Теорема Лебега о мажорированной сходимости.} Если $f_n\to f$ почти всюду
(или, в формулировке используемого конспекта на множестве конечной меры,
$f_n\to f$ по мере) и существует суммируемая функция $F$ такая, что
\[
    |f_n(x)|\leq F(x)\qquad\text{п.в. для всех }n,
\]
то
\[
    \int |f_n-f|\,d\mu\to0,
    \qquad
    \int f_n\,d\mu\to\int f\,d\mu.
\]

\emph{Теорема Беппо Леви} описывает монотонную сходимость: для
$0\leq f_n\uparrow f$ можно переходить к пределу под знаком интеграла,
\[
    \int f_n\,d\mu\uparrow\int f\,d\mu,
\]
с возможным значением $+\infty$. В частности, если интегралы $f_n$ равномерно
ограничены, предел суммируем.

\emph{Лемма Фату} для неотрицательных функций даёт оценку
\[
    \int \liminf_{n\to\infty}f_n\,d\mu
    \leq
    \liminf_{n\to\infty}\int f_n\,d\mu.
\]

\subsection{Связь с интегралом Римана}

На отрезке $[a,b]$ любая функция, интегрируемая по Риману, интегрируема и по
Лебегу, причём значения интегралов совпадают:
\[
    \int_a^b f(x)\,dx\bigg|_{\text{Риман}}
    =
    \int_a^b f(x)\,dx\bigg|_{\text{Лебег}}.
\]
Класс интегрируемых по Лебегу функций при этом существенно шире.

\subsection{Что важно сказать на экзамене}

Достаточно последовательно назвать $\sigma$-алгебру, меру и измеримую функцию,
после чего объяснить конструкцию интеграла Лебега: сначала для простых или
ограниченных измеримых функций, затем для неотрицательных и знакопеременных.
Нужно помнить критерий суммируемости $\int|f|d\mu<\infty$ и основные теоремы о
предельном переходе --- Лебега, Беппо Леви и Фату. Если функция интегрируема по
Риману, то она интегрируема по Лебегу и оба интеграла совпадают.

\newpage
